\section{Lezione 12} \hfill 24/03/2026
\subsubsection{Faggio a fustaia}
Gestione normale oltre le Alpi (non c'era la tradizione del ceduo e le aree sono superiori rispetto all'Italia).\\
La fustaia di faggio produce legname da opera (pavimenti, sedie, tavoli), che richiedono forte sollecitazioni.\\
Al momento il mercato è in crisi, data la competizione degli altri materiali.\\
Un buon utilizzo del faggio è il multistrato o il compensato.\\
La specie è definitiva, quindi in grado di rinnovare sotto se stessa, sciafila.\\
La gestione ottimale della faggeta sono i tagli successivi.\\
Tagli successivi: trattamento della fustaia coetania, consistente in più tagli, suddivisi in fondamentali e tagli accessori. I tagli fondamentali sono due: di sementazione e di sgombero. Quelli accessori sono i tagli di preparazione e quelli secondari. Il taglio di preparazione è un taglio da fare 20-30 anni prima del fine turno e consiste in un suer diradamento (in modo che le chiome si allarghino).\\
Una corretta gestione della fustaia non richiede il taglio....\\
Il primo taglio è quello di sementazione: la parola stessa indica l'obiettivo del taglio, ovvero avere il seme che germina e che si sviluppa. Il taglio di sementazione avviene in una annata di pasciona (ogni 5-6 anni), con una semipasciona ogni 2-3 anni.\\
Il taglio di sementazione consiste in una riduzione della copertura: arriva più luce al suolo e così la plantula ha più risorse di luce per svilupparsi. Non è necessario aprire grandi buche (solitamente 25-30\% di copertura). Su specie eliofile il taglio dev'essere superiore (50\%). Secondo il prof le buche per il faggio possono essere superiori.\\
Il vecchio popolamento, rimasto in campo, domina la rinnovazione: occorre liberare la rinnovazione. Si potrebbe eliminare tutto il vecchio, ma la rinnovazione è ancora molto piccola ed instabile. Viene così ridotta la densità del popolamento vecchio, mediante tagli secondari. Infine, si rimuove tutto il vecchio con un taglio di sgombero, quando la rinnovazione è affermata (2 metri).\\
Tra taglio di sementazione e di sgombero passano all'incirca 20 anni. Se dopo 20 anni non si riesce a fare il taglio di sgombero, vuol dire che è stato sbagliato qualcosa.\\
Questo sistema è tradizionale ed utilizzato dovunque. In ottica di biodiversità è limitante, perché si ottiene e rinnova solamente faggio (tutto della stessa età).\\
Sono stati ideati nuovi tagli (es. bonisky), dove c'è l'apertura di buche 2000-3000 metri, che permette l'ingresso di pioniere. Si incrementa così la biodiversità.\\
L'Italia è stato in passato molto ricco di cedui: ad un certo punto questi cedui non servivano più. Sono state fatte le conversioni dal ceduo di faggio a fustaie di faggio. Una fustaia di faggio produce sia legname da ardere che da opera, il ceduo solo da opera. Di contro, il ceduo ha dei turni molto più ridotti (20-30 anni ?). In Italia c'è una polverizzazione della proprietà privata.\\
Dagli anni 80 è iniziata la moda della conversione.\\
Le conversioni sono di moda tra i forestali.\\
Tradizionalmente in una conversione da ceduo a fustaia, ci sono diverse tipologie (nominato invecchiamento, svolto ad elevata età per non avere riscoppio dei polloni, 50-70 anni e dopo si fa un taglio di sementazione. questo sistema è stato abbandonato)
\begin{itemize}
    \item matricinatura intensiva: nel ceduo invecchiato si rilasciano molte matricine (il numero di polloni è di densità della fustaia ) , il periodo è di fustaia transitoria. si effettua il taglio di sementazione ..... i polloni hanno chiome sbandierate e perciò la qualità è piuttosto bassa.
    \item taglio di avviamento all'alto fusto: il ceduo viene avviato verso il governo a fustaia. Non conviene iniziarlo con il ceduo invecchiato (che ha le chiome sbandierate). Consiste in una serie di diradamenti, che riducono gradualmente il numero di polloni ad uno a ceppaia. Oltre che selezionare nel tempo le piante migliori, questo metodo permette di (ogni 10-15 anni) di fornire della legna.
\end{itemize}
Oggi non vengono più fatte le conversioni di faggi, perché c'è una discreta richiesta di legno di faggio da combustione.\\
Se l'area è servita da strade, se c'è mercato di legname da opera, se il faggio cresce bene (non è marginale): dove conviene fare le conversioni.\\
...\\
Oggi ci si pone il quesito del ritorno al ceduo, soprattutto per ceduare i vecchi cedui invecchiati.\\
Al ceduo invecchiato non bisogna pensare a turno di 120 anni, ma di 60-70 anni.\\
Provvigioni di 250-300 mq. Si potrebbe dare... ogni vent'anni.\\
\subsection{Abete bianco - Abies alba}
Regione esalpica o appenninica ?\\
L'abete bianco si riconosce rispetto all'abete rosso per:
\begin{itemize}
    \item coni verticali verso l'alto che maturano sulla pianta, si aprono ed il seme fuoriesce. Sul ramo rimane lo strobilo. L'abete rosso produce i coni verso l'alto ma quasi subito si abbassano
    \item l'abete rosso ha aghi con sezione romboidale, a punta, pungenti e di colore verde brillante. Gli aghi dell'abete bianco sono bianchi, finiscono arrotondati, non pungono, hanno la faccia superiore scura e la pagina inferiore bianca
    \item tutti e due portano ....     gli aghi di abete rosso si dispongono tutti attorno al ramo (formando quasi un manicotto)
    \item il cimale dell'abete bianco con l'invecchiamento tende a perdere dominanza, nell'abete rosso il cimale è sempre dominante;
    \item la corteccia nei giovani degli abeti rossi è bruna-rossastra, nell'abete bianco è grigia-biancastra. Nelle piante vecchie si confondono
    \item gli aghi di abete rosso hanno un odore di abete rosso, gli aghi di abete bianco hanno un odore di limone
    \item il legno dell'abete rosso ha canali resiniferi, il legno dell'abete bianco ha canali resiniferi nella corteccia
    \item il legno di abete rosso ha \dots
    \item l'abete rosso, in età adulta porta i rami orizzontalmente, mentre l'abete bianco porta i rami diagonalmente verso l'alto
    \item l'abete bianco è mediterranea, sciafila, definitiva (normalmente non è in grado di rinnovare sotto se stessa): quando è necessaria la rinnovazione, sotto l'abete bianco sono presenti altre specie (es. abete rosso o faggio) 
\end{itemize}
Il seme è abbastanza pesante, meno di quello del faggio. In natura è difficile avere abetine pure: solitamente sono opera dell'uomo (favorendo il bianco a discapito di altre specie o piantandolo). Abietino (piemontese o appenninico) o abieteto (veneto, friulano e trentita)\\
La specie è oceanica: climi umidi\\
L'apparato radicale è molto profondo, nelle formazioni dove si forma produce lo scheletro del terreno.\\
Il legno di abete bianco è stato per secoli odiato dai forestali: chi dice per la sua puzza (che comunque perde con l'essicazione) e chi dice che il legno ha troppo legno di reazione (dati i rami portati verso l'alto e non verso il basso come per l'abete rosso).\\
L'abete bianco è l'unica conifera nobile delle montagne appenniniche (tralasciando i pini). In appennino è stato molto diffuso e coltivato, per il proprio legname da opera.\\
Oggi l'abete bianco è stato rivalutato, potendo fare gli stessi usi dell'abete rosso.\\
È particolarmente apprezzato dai caprioli, perchè mangiato per gusto. In presenza di forti carichi di ungulati, la rinnovazione di abete bianco potrebbe scomparire.\\
La gestione delle abetine: se sono in purezza, difficilmente sono naturali (create e gestite). Per questo motivo occorre gestirle per tagli a buche e rinnovazione posticipata di abete bianco (per mantenere l'abetina) oppure le accompagniamo verso la formazione di boschi più misti, abbandonando la purezza.\\
Esistono diverse specie di abeti: es in Canada e nord-usa c'è la abies balsamea (specie definitiva). Nel bacino del mediterraneo ci sono diversi abeti bianchi. Il bianco è presente dalla Calabria, fino ai Volgi, passando per le Alpi. Sulle Madonie e Nebrodi c'è l abies nebrudensis ed è diversa dall'abete bianco. In Spagna c'è il abies pinsapo (aghi non perfettamente orizzontali, classico abete dei giardini). In Marocco c'è l Abies numidica. A cefalonica. in Turchia ci sono A cilica, A boriginensis e A nordmaniana (figli dello stesso abete iniziale).\\
L'abete bianco italiano delle alpi ha due origini genetiche: dopo le varie glaciazioni, è salito dall'appennino e dalla ex iugoslavia, trovandosi sulle alpi.\\
L'abete bianco ha avuto un'enorme crisi (deperimento) negli anni '80 (non si conoscono le cause). Oggi è in fortissima espansione. È consigliato mantenere e gestire il abies alba, come ``scheletro'' del bosco; inoltre, è utilizzabile per gli stessi fini dell'abete rosso.
\subsection{Pino silvestre}
I pini si dividono in due categorie: quelli con areali enormi e con areali piccolissimi. A volte i pini si ibridano, soprattutto nelle zone di contatto. Il pino silvestre ha l'areale più grande al mondo. Ad occidente sono i pirenei, a sud sono le Alpi, a Nord Norvegia, Svezia Finlandia, ad est la penisola di kamtchaka.\\
Commercialmente, in Italia, il legno viene venduto come pino di Scozia, di Russia, di Svezia. In inglese, si chiama Schotch pine.\\
Dal punto di vista ecologico, l'areale ampio permette di crescere in condizioni stazionali molto diverse.\\
COme tutti i pini è pioniera primaria (fortemente eliofila, apparato radicale fittonante, seme leggero, vita intorno ai 200-300 anni), ambiguitaria per il substrato (acidi, basici, sciolti, roccia viva o sabbia, torbosi),  